\documentclass[11pt,a4paper]{article}

% ─── Packages ─────────────────────────────────────────────────────
\usepackage[utf8]{inputenc}
\usepackage[T1]{fontenc}
\usepackage[margin=2.5cm]{geometry}
\usepackage{graphicx}
\usepackage{booktabs}
\usepackage{array}
\usepackage{amsmath}
\usepackage{hyperref}
\usepackage{xcolor}
\usepackage{fancyhdr}
\usepackage{float}
\usepackage{caption}
\usepackage{subcaption}
\usepackage{parskip}

% ─── Colors ───────────────────────────────────────────────────────
\definecolor{darkblue}{RGB}{0,51,102}
\definecolor{accentblue}{RGB}{30,100,180}

\hypersetup{
    colorlinks=true,
    linkcolor=darkblue,
    urlcolor=accentblue,
    citecolor=darkblue
}

% ─── Headers / Footers ────────────────────────────────────────────
\pagestyle{fancy}
\fancyhf{}
\rhead{\textcolor{darkblue}{\small SEDS 537 -- Machine Learning Midterm}}
\lhead{\textcolor{darkblue}{\small 323011020}}
\cfoot{\thepage}
\renewcommand{\headrulewidth}{0.4pt}

% ─── Document ─────────────────────────────────────────────────────
\begin{document}

% =================================================================
% TITLE PAGE
% =================================================================
\begin{titlepage}
    \centering
    \vspace*{1.5cm}

    {\Large \textbf{SEDS 537 -- Machine Learning} \par}
    \vspace{0.4cm}
    {\large Take-Home Midterm Examination \par}
    \vspace{0.2cm}
    {\small Instructor: Prof.\ Dr.\ Aytug Onan \par}

    \vspace{1.5cm}
    \rule{\linewidth}{1.5pt}
    \vspace{0.3cm}

    {\LARGE \textbf{Midterm Report} \par}
    \vspace{0.2cm}
    {\Large Five Machine Learning Tasks \par}
    \vspace{0.3cm}
    {\large Regression $\cdot$ Classification $\cdot$ Dimensionality Reduction \par}
    {\large Clustering $\cdot$ Neural Networks \par}

    \vspace{0.3cm}
    \rule{\linewidth}{1.5pt}

    \vspace{1.5cm}
    {\large \textbf{Student ID:} 323011020 \par}
    \vspace{0.3cm}
    {\large \textbf{Date:} April 2026 \par}

    \vfill
    {\small Department of Software Engineering \\
    Izmir Institute of Technology}
\end{titlepage}

% =================================================================
% TABLE OF CONTENTS
% =================================================================
\tableofcontents
\newpage

% =================================================================
% INTRODUCTION
% =================================================================
\section*{Introduction}
\addcontentsline{toc}{section}{Introduction}

This report presents the solutions for the five questions of the SEDS~537 Machine Learning take-home midterm.
Each question covers a different topic in machine learning:

\begin{itemize}
    \item \textbf{Question 1:} Regression on California Housing data.
    \item \textbf{Question 2:} Classification with class imbalance (Credit Card Fraud).
    \item \textbf{Question 3:} Dimensionality reduction with PCA and t-SNE on MNIST.
    \item \textbf{Question 4:} Clustering on the Wholesale Customer dataset.
    \item \textbf{Question 5:} Neural networks (MLP vs CNN) on Fashion-MNIST.
\end{itemize}

All experiments are reproducible. We used \texttt{random\_state=42} everywhere. The code is provided as Jupyter notebooks (Q1--Q5) and is publicly available at \url{https://github.com/yasntrk/experimental-ml-projects}. All figures shown in this report are generated from those notebooks.

\newpage

% =================================================================
% QUESTION 1 — REGRESSION
% =================================================================
\section{Question 1: Regression -- Predicting House Prices}

\subsection{Objective and Dataset}

The goal is to compare regression models on the \textbf{California Housing} dataset.
The dataset has 20{,}640 samples and 8 numerical features. The target is the median house value (in \$100{,}000s).

\subsection{Exploratory Data Analysis}

We first looked at feature distributions, correlations, and scatter plots.

\begin{figure}[H]
    \centering
    \includegraphics[width=0.95\textwidth]{q1_histograms.png}
    \caption{Histograms of all features in the California Housing dataset.}
    \label{fig:q1_hist}
\end{figure}

\begin{figure}[H]
    \centering
    \includegraphics[width=0.7\textwidth]{q1_correlation_matrix.png}
    \caption{Correlation matrix. Median income (\texttt{MedInc}) has the strongest positive correlation with the target.}
    \label{fig:q1_corr}
\end{figure}

\begin{figure}[H]
    \centering
    \includegraphics[width=\textwidth]{q1_scatter_plots.png}
    \caption{Scatter plots of the four most correlated features against the target.}
    \label{fig:q1_scatter}
\end{figure}

\textbf{Observations:}
\begin{itemize}
    \item \texttt{MedInc} is highly correlated with the target ($r \approx 0.69$).
    \item Some features (\texttt{AveRooms}, \texttt{Population}) are very right-skewed.
    \item There is a clear cap at \texttt{MedHouseVal}~$=5.0$ which appears as a horizontal band in the scatter plots.
\end{itemize}

\subsection{Models and Evaluation}

We compared three regression models: Linear Regression, Ridge Regression with cross-validated alpha, and Random Forest. We evaluated them with three settings: original features, scaled features, and polynomial (degree 2) features.

\begin{table}[H]
\centering
\caption{Q1 Results -- All Models and Settings}
\label{tab:q1_results}
\begin{tabular}{lccc}
\toprule
\textbf{Model} & \textbf{RMSE} & \textbf{MAE} & \textbf{R²} \\
\midrule
Linear Regression          & 0.7456 & 0.5332 & 0.5758 \\
Ridge Regression (CV)      & 0.7451 & 0.5332 & 0.5763 \\
Random Forest              & 0.5057 & 0.3276 & 0.8049 \\
\midrule
Linear Regression (Scaled) & 0.7456 & 0.5332 & 0.5758 \\
Ridge Regression (Scaled)  & 0.7456 & 0.5332 & 0.5758 \\
Random Forest (Scaled)     & 0.5055 & 0.3276 & 0.8050 \\
\midrule
Linear Regression (Poly)   & 0.6814 & 0.4670 & 0.6457 \\
Ridge Regression (Poly)    & 0.6787 & 0.4819 & 0.6485 \\
Random Forest (Poly)       & 0.5126 & 0.3332 & 0.7995 \\
\bottomrule
\end{tabular}
\end{table}

\subsection{Impact of Scaling}

Standard scaling makes all features have mean 0 and standard deviation 1.
\begin{itemize}
    \item For \textbf{Linear Regression}, scaling does not change predictions; only coefficient sizes change.
    \item For \textbf{Ridge Regression}, scaling is important because the L2 penalty treats all features equally only when they are on the same scale.
    \item For \textbf{Random Forest}, scaling has no effect because trees are based on splits.
\end{itemize}

\subsection{Polynomial Features}

We added polynomial features of degree 2 (44 features from 8 original).
This improved Linear and Ridge Regression a lot (R²: 0.58 $\to$ 0.65). For Random Forest, polynomial features did not help, because the model can already capture non-linear patterns.

\subsection{Residual Plots}

\begin{figure}[H]
    \centering
    \includegraphics[width=\textwidth]{q1_residual_plots.png}
    \caption{Residual plots (Fitted values vs Residuals) for the three models.}
    \label{fig:q1_residuals}
\end{figure}

\textbf{Analysis:}
\begin{itemize}
    \item Linear and Ridge residuals show a funnel shape (heteroscedasticity). Variance increases with predicted value.
    \item There is a horizontal band at the top because the target is capped at 5.0.
    \item Random Forest residuals are more uniform around zero, which means the model fits better.
\end{itemize}

\subsection{Conclusion: Interpretability vs Predictive Performance}

\begin{table}[H]
\centering
\caption{Trade-off between Interpretability and Performance}
\begin{tabular}{lll}
\toprule
\textbf{Model} & \textbf{Interpretability} & \textbf{Performance} \\
\midrule
Linear Regression & High & Moderate \\
Ridge Regression  & High & Moderate \\
Random Forest     & Low (black box) & High \\
\bottomrule
\end{tabular}
\end{table}

Random Forest gives the best results (R²~$\approx 0.80$) but is hard to interpret.
Linear models are easy to explain but cannot capture non-linear patterns.
Polynomial features bring linear models closer to Random Forest performance with the cost of more features.

\newpage

% =================================================================
% QUESTION 2 — CLASSIFICATION (IMBALANCE)
% =================================================================
\section{Question 2: Classification -- Handling Class Imbalance}

\subsection{Dataset and Class Distribution}

We used the \textbf{Credit Card Fraud Detection} dataset from Kaggle.
It has 284{,}807 transactions and only 492 fraud cases (about 0.17\%).

\begin{table}[H]
\centering
\caption{Class Distribution}
\begin{tabular}{lrr}
\toprule
\textbf{Class} & \textbf{Count} & \textbf{Percentage} \\
\midrule
Normal (0) & 284{,}315 & 99.827\% \\
Fraud  (1) &     492   &  0.173\% \\
\bottomrule
\end{tabular}
\end{table}

\begin{figure}[H]
    \centering
    \includegraphics[width=0.55\textwidth]{q2_class_distribution.png}
    \caption{Class Distribution: Normal vs Fraud Transactions.}
    \label{fig:q2_classdist}
\end{figure}

The dataset is very imbalanced. A model that always predicts ``Normal'' would have 99.83\% accuracy, but it would not catch any fraud. So we need better metrics (precision, recall, F1, ROC-AUC) and resampling techniques.

\subsection{Preprocessing}

\begin{itemize}
    \item Scaled \texttt{Amount} and \texttt{Time} with \texttt{StandardScaler} (V1--V28 are already PCA-transformed).
    \item Train/Test split: 80\%/20\% with \texttt{stratify=y}.
\end{itemize}

\subsection{Classifiers}

We trained three classifiers: \textbf{Logistic Regression}, \textbf{Random Forest}, and \textbf{XGBoost}.

\subsection{Resampling Techniques}

\begin{itemize}
    \item \textbf{SMOTE:} creates synthetic fraud samples by interpolating between existing fraud points.
    \item \textbf{Random undersampling:} removes random normal samples until the classes are balanced.
\end{itemize}

\subsection{Results}

\begin{table}[H]
\centering
\caption{Q2 Results -- All Models and Resampling Strategies (Fraud Class)}
\label{tab:q2_results}
\begin{tabular}{llcccc}
\toprule
\textbf{Resampling} & \textbf{Model} & \textbf{Precision} & \textbf{Recall} & \textbf{F1} & \textbf{ROC-AUC} \\
\midrule
None        & Logistic Regression & 0.8713 & 0.6327 & 0.7333 & 0.9691 \\
None        & Random Forest       & 0.9524 & 0.8163 & 0.8791 & 0.9771 \\
None        & XGBoost             & 0.9302 & 0.8163 & 0.8696 & 0.9814 \\
\midrule
SMOTE       & Logistic Regression & 0.0711 & 0.9184 & 0.1317 & 0.9724 \\
SMOTE       & Random Forest       & 0.8947 & 0.8673 & 0.8808 & 0.9782 \\
SMOTE       & XGBoost             & 0.8148 & 0.8980 & 0.8545 & 0.9813 \\
\midrule
Undersample & Logistic Regression & 0.0524 & 0.9388 & 0.0993 & 0.9641 \\
Undersample & Random Forest       & 0.0679 & 0.9184 & 0.1268 & 0.9718 \\
Undersample & XGBoost             & 0.0812 & 0.9082 & 0.1491 & 0.9756 \\
\bottomrule
\end{tabular}
\end{table}

\subsection{Confusion Matrices and ROC Curves}

\begin{figure}[H]
    \centering
    \includegraphics[width=\textwidth]{q2_confusion_original.png}
    \caption{Confusion matrices without resampling.}
    \label{fig:q2_cm}
\end{figure}

\begin{figure}[H]
    \centering
    \includegraphics[width=0.75\textwidth]{q2_roc_curves.png}
    \caption{ROC curves for XGBoost under different resampling strategies.}
    \label{fig:q2_roc}
\end{figure}

\subsection{Precision-Recall Trade-off}

In fraud detection:
\begin{itemize}
    \item \textbf{False Negative (FN):} missed fraud. The bank loses money. \emph{Very costly}.
    \item \textbf{False Positive (FP):} a normal transaction marked as fraud. \emph{Less costly} -- only inconvenience.
\end{itemize}

So \textbf{Recall is more important than Precision}. We prefer to catch as much fraud as possible, even with some false alarms.

\subsection{Conclusion}

\begin{itemize}
    \item XGBoost has the highest ROC-AUC ($\approx 0.98$).
    \item SMOTE is better than random undersampling: it gives higher recall while keeping precision reasonable. Undersampling drops precision very low.
    \item \textbf{Random Forest with SMOTE} (F1 = 0.8808) gives the best balance for this task.
\end{itemize}

\newpage

% =================================================================
% QUESTION 3 — DIMENSIONALITY REDUCTION
% =================================================================
\section{Question 3: Dimensionality Reduction and Visualisation}

\subsection{Dataset}

We used the \textbf{MNIST} dataset (70{,}000 handwritten digit images, 28$\times$28 pixels = 784 features).
For visualisation and t-SNE, we used a random subset of 8{,}000--10{,}000 samples (t-SNE is slow for large data).

\subsection{PCA: 2D and 50D}

We reduced data to 2D and 50D with PCA.

\begin{figure}[H]
    \centering
    \includegraphics[width=\textwidth]{q3_pca_variance.png}
    \caption{Left: explained variance ratio for PC1--PC10. Right: cumulative explained variance for 50 PCs.}
    \label{fig:q3_pca_var}
\end{figure}

\textbf{Explained variance for first 10 components:}

\begin{table}[H]
\centering
\caption{PCA Explained Variance Ratio (First 10 Components)}
\begin{tabular}{lcccccccccc}
\toprule
\textbf{PC} & 1 & 2 & 3 & 4 & 5 & 6 & 7 & 8 & 9 & 10 \\
\midrule
\textbf{Ratio} & 0.097 & 0.071 & 0.062 & 0.054 & 0.049 & 0.043 & 0.033 & 0.029 & 0.028 & 0.024 \\
\bottomrule
\end{tabular}
\end{table}

The first 10 components explain about 49\% of variance. The first 50 components explain about 82\%.

\subsection{t-SNE: 2D Embeddings}

t-SNE was applied with two perplexity values: 30 and 50.

\begin{figure}[H]
    \centering
    \includegraphics[width=\textwidth]{q3_embeddings.png}
    \caption{2D embeddings: PCA (left), t-SNE perplexity~30 (middle), t-SNE perplexity~50 (right).}
    \label{fig:q3_emb}
\end{figure}

\textbf{Observations:}
\begin{itemize}
    \item PCA 2D projection shows overlapping classes -- it cannot separate digits well in 2D.
    \item t-SNE produces clear clusters for each digit. The clusters are well separated.
    \item Higher perplexity (50) makes clusters bigger and slightly closer to each other.
\end{itemize}

\subsection{k-NN Classification Comparison}

We trained k-NN (k=5) on three feature spaces and used 5-fold cross-validation.

\begin{table}[H]
\centering
\caption{Q3 -- 5-Fold CV Accuracy of k-NN}
\label{tab:q3_knn}
\begin{tabular}{lc}
\toprule
\textbf{Feature Space} & \textbf{Accuracy} \\
\midrule
Original (784D)        & $\approx 0.95$ \\
PCA 50D                & $\approx 0.95$ \\
t-SNE 2D               & $\approx 0.96$ \\
\bottomrule
\end{tabular}
\end{table}

\subsection{PCA vs t-SNE Discussion}

\begin{table}[H]
\centering
\caption{PCA vs t-SNE Comparison}
\begin{tabular}{p{3.5cm}p{5.5cm}p{5.5cm}}
\toprule
\textbf{Aspect} & \textbf{PCA} & \textbf{t-SNE} \\
\midrule
Type           & Linear, deterministic       & Non-linear, stochastic \\
Speed          & Fast                        & Slow \\
Visualisation  & Overlapping clusters        & Well-separated clusters \\
Downstream use & Good (preserves global structure) & Limited (only local) \\
New data       & Yes (transform method)      & No (must refit) \\
\bottomrule
\end{tabular}
\end{table}

\textbf{Conclusion:} PCA is best for preprocessing before classification. t-SNE is best for visualization and exploratory analysis.

\newpage

% =================================================================
% QUESTION 4 — CLUSTERING
% =================================================================
\section{Question 4: Clustering -- Unsupervised Learning}

\subsection{Dataset}

We used the \textbf{Wholesale Customer} dataset from UCI (440 customers, 6 spending features: Fresh, Milk, Grocery, Frozen, Detergents\_Paper, Delicassen).
Features were scaled with \texttt{StandardScaler}.

\subsection{K-Means: Choosing k}

\begin{figure}[H]
    \centering
    \includegraphics[width=\textwidth]{q4_elbow_silhouette.png}
    \caption{Left: Elbow method (Inertia vs k). Right: Silhouette score vs k.}
    \label{fig:q4_elbow}
\end{figure}

The silhouette score is highest at $k = 2$. So we used $k = 2$ for K-Means.

\subsection{Agglomerative Clustering (Ward Linkage)}

\begin{figure}[H]
    \centering
    \includegraphics[width=0.85\textwidth]{q4_dendrogram.png}
    \caption{Hierarchical clustering dendrogram (Ward linkage).}
    \label{fig:q4_dendro}
\end{figure}

The dendrogram supports 2 large clusters as well.

\subsection{DBSCAN: Tuning eps}

\begin{figure}[H]
    \centering
    \includegraphics[width=0.7\textwidth]{q4_kdistance.png}
    \caption{k-distance graph (k=5) for DBSCAN \texttt{eps} selection.}
    \label{fig:q4_kdist}
\end{figure}

We chose \texttt{eps} from the elbow of this curve.

\subsection{Evaluation Metrics}

\begin{table}[H]
\centering
\caption{Q4 -- Internal Cluster Evaluation Metrics}
\label{tab:q4_metrics}
\begin{tabular}{lcc}
\toprule
\textbf{Algorithm}   & \textbf{Silhouette} (higher is better) & \textbf{Davies-Bouldin} (lower is better) \\
\midrule
K-Means (k=2)        & $\approx 0.50$  & $\approx 0.81$ \\
Agglomerative (k=2)  & $\approx 0.49$  & $\approx 0.85$ \\
DBSCAN               & varies          & varies \\
\bottomrule
\end{tabular}
\end{table}

\subsection{Cluster Visualisation}

\begin{figure}[H]
    \centering
    \includegraphics[width=\textwidth]{q4_cluster_visualization.png}
    \caption{Clusters projected to 2D using PCA, for each algorithm.}
    \label{fig:q4_clusters}
\end{figure}

\subsection{Cluster Interpretation}

\begin{figure}[H]
    \centering
    \includegraphics[width=0.85\textwidth]{q4_cluster_profiles.png}
    \caption{K-Means cluster profiles: normalized average spending per category.}
    \label{fig:q4_profiles}
\end{figure}

\textbf{Cluster characteristics:}
\begin{itemize}
    \item \textbf{Cluster 0:} Customers with moderate spending in all categories. Likely small retailers or restaurants (high \texttt{Fresh}).
    \item \textbf{Cluster 1:} Customers with higher spending on \texttt{Milk}, \texttt{Grocery}, and \texttt{Detergents\_Paper}. Likely supermarkets or hotels.
\end{itemize}

\subsection{DBSCAN Sensitivity}

DBSCAN is very sensitive to its parameters:
\begin{itemize}
    \item Very small \texttt{eps}: most points become noise; few or no clusters.
    \item Very large \texttt{eps}: all points merge into one cluster.
    \item DBSCAN can find arbitrary-shaped clusters and handle outliers, but it struggles when clusters have different densities.
\end{itemize}

\subsection{Algorithm Comparison}

\begin{table}[H]
\centering
\caption{Clustering Algorithms Compared}
\begin{tabular}{lll}
\toprule
\textbf{Algorithm} & \textbf{Cluster Shape} & \textbf{Notes} \\
\midrule
K-Means        & Spherical & Fast; needs k; sensitive to outliers \\
Agglomerative  & Hierarchical & Dendrogram useful; expensive memory \\
DBSCAN         & Arbitrary & Finds noise; sensitive to eps and density \\
\bottomrule
\end{tabular}
\end{table}

\newpage

% =================================================================
% QUESTION 5 — NEURAL NETWORKS
% =================================================================
\section{Question 5: Neural Networks -- MLP vs CNN}

\subsection{Dataset}

We used \textbf{Fashion-MNIST}: 70{,}000 grayscale images of 10 clothing classes (28$\times$28 pixels).
Train/Val/Test split: 51{,}000 / 9{,}000 / 10{,}000 samples.

\subsection{Models}

\textbf{MLP (Multilayer Perceptron):}
\begin{itemize}
    \item Input flattened to 784 values.
    \item Hidden layer 1: 256 neurons, ReLU, Dropout 0.3.
    \item Hidden layer 2: 128 neurons, ReLU, Dropout 0.3.
    \item Output: 10 classes.
\end{itemize}

\textbf{CNN (Convolutional Neural Network):}
\begin{itemize}
    \item Conv2D (32 filters, 3$\times$3) + ReLU + MaxPool 2$\times$2.
    \item Conv2D (64 filters, 3$\times$3) + ReLU + MaxPool 2$\times$2.
    \item Flatten + Dense (128) + ReLU + Dropout 0.3.
    \item Output: 10 classes.
\end{itemize}

Both models were trained with Adam optimizer (lr=0.001) and cross-entropy loss. Early stopping used validation loss with patience~4.

\subsection{Training Curves}

\begin{figure}[H]
    \centering
    \includegraphics[width=\textwidth]{q5_training_curves.png}
    \caption{Training and validation curves for MLP (left) and CNN (right).}
    \label{fig:q5_curves}
\end{figure}

\textbf{Overfitting analysis:}
\begin{itemize}
    \item Both models show some overfitting after a few epochs (training loss keeps going down, validation loss increases).
    \item Early stopping helps to keep the best model.
    \item Dropout layers and regularization help to reduce overfitting.
\end{itemize}

\subsection{Test Set Performance}

\begin{table}[H]
\centering
\caption{Q5 -- Test Performance}
\label{tab:q5_results}
\begin{tabular}{lcc}
\toprule
\textbf{Model} & \textbf{Accuracy} & \textbf{Macro-F1} \\
\midrule
MLP & $\approx 0.88$ & $\approx 0.88$ \\
CNN & $\approx 0.91$ & $\approx 0.91$ \\
\bottomrule
\end{tabular}
\end{table}

\subsection{Confusion Matrices}

\begin{figure}[H]
    \centering
    \includegraphics[width=\textwidth]{q5_confusion_matrices.png}
    \caption{Confusion matrices on the test set for MLP and CNN.}
    \label{fig:q5_cm}
\end{figure}

The most confused pairs are: \emph{Shirt} ↔ \emph{T-shirt/top} ↔ \emph{Coat} ↔ \emph{Pullover}. These classes share similar shapes (sleeves, collars).

\subsection{Misclassified Examples}

\begin{figure}[H]
    \centering
    \includegraphics[width=\textwidth]{q5_misclassified_mlp.png}
    \caption{Five MLP misclassified samples.}
    \label{fig:q5_mis_mlp}
\end{figure}

\begin{figure}[H]
    \centering
    \includegraphics[width=\textwidth]{q5_misclassified_cnn.png}
    \caption{Five CNN misclassified samples.}
    \label{fig:q5_mis_cnn}
\end{figure}

\textbf{Possible reasons for misclassification:}
\begin{itemize}
    \item Low resolution (28$\times$28) makes it hard to see details.
    \item Some classes are visually very similar (Shirt vs T-shirt vs Pullover).
    \item Sandals and sneakers can look similar in low-resolution silhouettes.
\end{itemize}

\subsection{Why CNN Performs Better}

\begin{table}[H]
\centering
\caption{MLP vs CNN Comparison}
\begin{tabular}{p{4cm}p{5cm}p{5cm}}
\toprule
\textbf{Aspect} & \textbf{MLP} & \textbf{CNN} \\
\midrule
Input handling      & Flattened (784,) -- spatial info lost & 2D (1, 28, 28) -- spatial info preserved \\
Feature detection   & Position-specific & Translation-invariant (filters slide over image) \\
Parameter sharing   & No -- huge dense layers & Yes -- weights shared across image \\
Hierarchical features & Limited & Strong (edges $\to$ shapes $\to$ objects) \\
\bottomrule
\end{tabular}
\end{table}

\textbf{Conclusion:} CNN performs better on images because:
\begin{enumerate}
    \item Convolutional layers detect local patterns (edges, textures) regardless of where they appear in the image.
    \item Weight sharing reduces parameters and prevents overfitting.
    \item Pooling makes the model robust to small shifts.
    \item Hierarchical feature learning matches the structure of natural images.
\end{enumerate}

\newpage

% =================================================================
% FINAL CONCLUSION
% =================================================================
\section{Overall Conclusion}

This midterm covered five core topics in machine learning. The main takeaways are:

\begin{enumerate}
    \item \textbf{Q1 (Regression):} Random Forest beats linear models on California Housing. Polynomial features help linear models but increase complexity.
    \item \textbf{Q2 (Classification):} For imbalanced data, accuracy is misleading. SMOTE + Random Forest gives the best precision/recall balance for fraud detection.
    \item \textbf{Q3 (Dim. Reduction):} PCA preserves global structure (good for downstream tasks); t-SNE produces beautiful visualizations of local structure.
    \item \textbf{Q4 (Clustering):} K-Means with $k=2$ separates wholesale customers well into ``small retailers'' and ``supermarkets''. DBSCAN is powerful but very sensitive to parameters.
    \item \textbf{Q5 (Neural Networks):} CNNs clearly outperform MLPs on image data, thanks to translation invariance and hierarchical feature learning.
\end{enumerate}

The general lesson is: \emph{the right model depends on the data}. Linear models are good when relationships are simple. Tree ensembles work well on tabular data. Neural networks excel on unstructured data like images.

\vspace{1cm}
\noindent\rule{\linewidth}{0.5pt}
\begin{center}
    \small \textit{SEDS 537 -- Machine Learning Take-Home Midterm} \\
    \textit{Student ID: 323011020 | April 2026}
\end{center}

\end{document}
